% Part: counterfactuals
% Chapter: introduction
% Section: material-conditional

\documentclass[../../../include/open-logic-section]{subfiles}

\begin{document}

\olfileid[pt]{cnt}{int}{mat}

\olsection{O condicional material}

Em sua forma mais simples em inglês, um condicional é uma sentença da
forma ``Se \dots, então \dots'', em que os \dots{} são, eles próprios,
sentenças, como ``Se o mordomo fez isso, então o jardineiro é
inocente''. Em cursos introdutórios de lógica, aprendemos a simbolizar
condicionais usando o conectivo $\lif$: simbolizamos as partes
indicadas por \dots, p.ex., pelas !!{formula}s $!A$ e~$!B$,
e o condicional inteiro é simbolizado por $!A \lif !B$.

O conectivo $\lif$ é \emph{verofuncional}, O.s.\@, o valor de
verdade---$\True$ ou $\False$---de $!A \lif !B$ é determinado pelos valores de
verdade de $!A$ e~$!B$: $!A \lif !B$ é verdadeiro se, e somente se, $!A$~é falso ou
$!B$~é verdadeiro, e é falso nos demais casos. Relativamente a uma
atribuição de valores de verdade~$\pAssign v$, definimos $\pSat{v}{!A \lif !B}$ se, e somente se,
$\pSat/{v}{!A}$ ou $\pSat{v}{!B}$. O conectivo $\lif$ com essa
semântica é chamado de \emph{condicional material}.

Essa definição produz vários fatos lógicos elementares. Antes
de tudo, o teorema da dedução vale para o condicional material:
\begin{align}
  \text{Se } \Gamma, !A \Entails !B \text{, então } \Gamma \Entails !A \lif !B
\end{align}
Ele é verofuncional: $!A \lif !B$ e $\lnot !A \lor !B$ são equivalentes:
\begin{align}
  !A \lif !B & \Entails \lnot !A \lor !B\\
  \lnot !A \lor !B & \Entails !A \lif !B
  \intertext{Um condicional material é implicado por seu consequente e
    pela negação de seu antecedente:}
  !B & \Entails !A \lif !B\\
  \lnot !A & \Entails !A \lif !B
  \intertext{Um condicional material falso equivale à
    conjunção de seu antecedente com a negação de seu consequente:
    se $!A \lif !B$ é falso, $!A \land \lnot !B$ é verdadeiro, e vice-versa:}
  \lnot(!A \lif !B) & \Entails !A \land \lnot !B\\
  !A \land \lnot !B & \Entails \lnot(!A \lif !B)
  \intertext{O condicional material admite o modus ponens:}
  !A, !A \lif !B & \Entails !B
  \intertext{O condicional material aglomera:}
  !A \lif !B,  !A \lif !C & \Entails !A \lif (!B \land !C)
  \intertext{Sempre podemos fortalecer o antecedente, O.s.\@, o
    condicional é \emph{monótono}:}
  !A \lif !B & \Entails (!A \land !C) \lif !B
  \intertext{O condicional material é transitivo, O.s.\@, a regra da
    cadeia é válida:}
  !A \lif !B, !B \lif !C & \Entails !A \lif !C
  \intertext{O condicional material equivale à sua
    contrapositiva:}
  !A \lif !B & \Entails \lnot !B \lif \lnot !A\\
  \lnot !B \lif \lnot !A & \Entails !A \lif !B
\end{align}

Todas essas são inferências úteis e não problemáticas no raciocínio
matemático. Contudo, a literatura filosófica e linguística está
repleta de supostos contraexemplos às inferências equivalentes em
contextos não matemáticos. Eles sugerem que o condicional
material~$\lif$ não é---ou pelo menos nem sempre é---o conectivo
apropriado para simbolizar enunciados em inglês da forma ``se \dots,
então \dots''.

\end{document}
